\subsection{A complete pointed input census and first character fields}

\begin{claim}{B.81}{Proved}{Pointed quartic representatives and the first possible character fields}
There are exactly $3001$ pointed smooth plane-quartic isomorphism
classes over $\mathbf F _3$, and exactly $94283$ over $\mathbf F _5$.
Here a point is part of the object and is rational over the stated
base field; isomorphisms preserve it. Complete equations, point
counts over the first three extensions, and Weil polynomials are
retained for every class.

For each such pointed curve let $m_\ell$ be the least positive
integer $m$ for which a nonzero character
$J_C(\mathbf F _{q^m})\twoheadrightarrow\mathbf Z/\ell$ exists.
For the smallest admissible odd prime, namely $\ell=5$ when $q=3$
and $\ell=3$ when $q=5$, the exact distributions are
\[
\begin{array}{c|rrrrrr}
q=3:\ m_5&1&2&3&6&8&24\\ \hline
\text{classes}&650&515&49&86&327&350\\[2pt]
m_5&31&56&62&72&168&504\\ \hline
\text{classes}&255&47&261&79&73&309
\end{array}
\]
and
\[
\begin{array}{c|rrrrr}
q=5:\ m_3&1&4&8&13&56\\ \hline
\text{classes}&41623&5163&19699&14400&13398.
\end{array}
\]
These are first fields of existence of a nonzero constant character,
not first fields of a class-number witness.
\end{claim}

\begin{proof}
A smooth plane quartic is canonically embedded, so its abstract
isomorphisms are exactly the projective linear transformations of
its equation. Send the marked point to $P=(0:0:1)$ and its tangent
to $Y=0$, then normalize the coefficient of $YZ^3$ to one. Its
unique general form is
\[
\begin{split}
 F={}&YZ^3+(aX^2+bXY+cY^2)Z^2\\
     &+(dX^3+eX^2Y+fXY^2+gY^3)Z+Q_4(X,Y).
\end{split}                                                   \tag{FC.10}
\]
The stabilizer is
\[
 (X,Y,Z)\longmapsto(AX+BY,CY,DX+EY+Z),\qquad A,C\ne0,
                                                               \tag{FC.11}
\]
followed by division of the polynomial by $C$. Over $\mathbf F _3$
this group has order $108$. Enumerate all $3^{12}=531441$ coefficient
vectors and mark their full orbits under (FC.11).

For $\mathbf F _5$ the following three slices suffice:
\[
\begin{array}{c|c|c|c}
\text{tangent contact}&\text{fixed coefficients}&\text{free}&w\\ \hline
2&a=1,\ b=c=d=0&8&2\\
3&a=0,\ d=1,\ b=c=e=0&7&3\\
4&a=d=0,\ u=1,\ b=c=v=0&6&4,
\end{array}
\]
where $u,v$ are the coefficients of $X^4,X^3Y$ in $Q_4$.
Here is a proof of completeness and uniqueness up to the stated
residual group. The transformed first coefficients are
\[
\begin{aligned}
a'&=aA^2/C,&b'&=3D+2aAB/C+bA,\\
c'&=3E+aB^2/C+bB+cC,&d'&=2aA^2D/C+dA^3/C.
\end{aligned}
\]
If $a\ne0$, choose $C=aA^2$, $D=-dA/(2a)$,
$B=-A(3D+bA)/2$, then the unique $E$ for which $c'=0$.
If $a=0,d\ne0$, choose $C=dA^3$, $D=-bA/3$, then the unique
$B$ making $e'=3D^2+2bAD+3B/A+eA^2$ zero, and again the
unique $E$ making $c'=0$.
If $a=d=0$, smoothness forces $u\ne0$, since otherwise the tangent
line is a component. Choose $C=uA^4$, $D=-bA/3$, then the unique
$B$ making $v'=D^3+bAD^2+eA^2D+4B/A+vA^3$ zero and the unique
$E$ making $c'=0$. All denominators are nonzero in characteristic
five. Conversely these coefficient equations show that a map between
two forms of the same slice has $C=A^w$ and $D=B=E=0$.
Its action on the coefficient of $X^iY^jZ^{4-i-j}$ is therefore
multiplication by $A^{i+wj-w}$. The residual group has order four.
This reduces the enumeration to $5^8+5^7+5^6=484375$ forms under
four diagonal transformations. The contact orders distinguish the
three slices intrinsically.

Smoothness is tested over the algebraic closure, not merely at
rational points. In three variables, multiply each of the three
cubic partial derivatives by every degree-four monomial and form
the resulting $45\times36$ coefficient matrix in degree seven.
The curve is smooth exactly when this matrix has rank $36$.
Indeed Euler's identity has invertible coefficient four in both
characteristics. A common nonzero zero of the partials makes degree
seven nonzero in their homogeneous quotient. In the absence of such
a zero the three cubic partials form a regular sequence in the
three-variable polynomial ring. Its quotient has Hilbert series
$(1+T+T^2)^3$ and hence no degree-seven part. These arguments prove
both directions of the rank test.

The complete enumeration gives $3001$ smooth orbits, with total
orbit size $315360$, over $\mathbf F _3$. Over $\mathbf F _5$ the
three slices give $75418,15626,3239$ smooth orbits. Their sum is
$94283$. An orbit of size $s$ in a slice corresponds to $500s$
forms in (FC.10), because the full stabilizer has order $2000$.
The weighted total is $187512000$. These identities are checked
by the program itself, and every smooth representative, its orbit
size and contact order is retained in the output tables.

For each representative, direct point enumeration gives
$N_i=\#C(\mathbf F _{q^i})$ for $i=1,2,3$. The implementation
uses the chart $Y=1$, in which (FC.10) is monic cubic in $Z$;
precomputed exact cubic-root counts over the six small fields
avoid repeated polynomial solving. The chart $Y=0$ is counted
separately. The complete genus-three Frobenius polynomial is
\[
 H(U)=U^6+a_1U^5+a_2U^4+a_3U^3+qa_2U^2+q^2a_1U+q^3,
\]
where, putting $s_i=q^i+1-N_i$, Newton's identities give
\[
 a_1=-s_1,\qquad
 2a_2=s_1^2-s_2,\qquad
 3a_3=-(s_3+a_1s_2+a_2s_1).
\]
These are integer identities, so the divisions are taken before
reduction modulo $\ell$.

A nonzero quotient to $\mathbf Z/\ell$ exists exactly when
$\ell$ divides the finite abelian group order
$\#J_C(\mathbf F _{q^m})$. Modulo $\ell$ this order is
$\prod_j(1-\lambda_j^m)$, where the $\lambda_j$ are the roots
of $H$ in an algebraic closure of $\mathbf F _\ell$. Therefore
\[
 m_\ell=\min_f\operatorname{ord}(U\bmod f),                  \tag{FC.12}
\]
with $f$ ranging over the distinct monic irreducible factors of
$H$ modulo $\ell$. This does not assume semisimplicity: the
singularity of $F^m-1$, or simply its determinant, is enough.
For a factor of degree $d$, the root order divides $\ell^d-1$;
exact modular powering, removing prime divisors from that integer,
computes the order. Every factor and order is retained. Applying
(FC.12) to the full tables gives the stated distributions.
The pointed Lang pullback realizes each resulting quotient with
its unique split fibre at $P$, so the character existence has its
usual geometric interpretation.
\end{proof}

\paragraph{Check record and boundary.}
Also checked in this writing for the census totals carried from
AS-A.72--AS-A.73, and checked only in this writing for the complete
first-character-field distributions: GPT-6 Astra Pro, 18 November
2026, finite-checks reader. The programs
\path{code/pointed_quartic_census.cpp},
\path{code/quartic_weil_census.cpp}, and
\path{code/quartic_character_fields.py} were written from the
mathematical descriptions, without access to the absent prior
census programs or their large data files. Their equations, all
point counts, Weil polynomials, factors and root orders are retained.
The effort setting is not separately exposed.

The large first degrees $56$ and $504$ show why a search through a
few constant extensions cannot decide B6. This statement does not
compute the required Cartier action on the new Jacobian parts at
those fields. It also concerns quartics with a point over the base
field; curves acquiring their first point only after extension
remain part of the uncompleted universal question.
